% Compile with pdfLaTeX (the bundled ByteSans fonts use pdfTeX mappings).
\documentclass[twocolumn]{paper}
% Jilin University theme overrides; paper.cls only renames the theme color.
% geometry also uses addtolist; preserve the class metadata helper.
\let\jlusavedaddtolist\addtolist
\geometry{a4paper,headheight=43pt}
\let\addtolist\jlusavedaddtolist
\DeclareTextFontCommand{\textbf}{\bfseries}
\paperlogo{figures/jlu.pdf}
\fancypagestyle{firststyle}{
  \fancyhf{}
  \fancyfoot[C]{\thepage}
  \fancyhead[L]{{\frontmedium\small\color{jlublue}JILIN UNIVERSITY}}
  \fancyhead[R]{%
    \ifdefempty{\paperlogofile}{}{\includegraphics[height=13.5mm]{\paperlogofile}}%
  }
}

% Add paper-specific packages and commands here.
\usepackage{amsmath,amssymb,amsthm}
\usepackage{enumitem}
\usepackage{tikz}
\graphicspath{{figures/}}
\setlist[itemize]{leftmargin=*,nosep}
\newtheorem{theorem}{Theorem}[section]
\newtheorem{lemma}[theorem]{Lemma}
\theoremstyle{definition}
\newtheorem{definition}[theorem]{Definition}
\theoremstyle{remark}
\newtheorem{remark}[theorem]{Remark}
\newcommand{\vect}[1]{\bm{#1}}
\newcommand{\R}{\mathbb{R}}

% Highlight box, matching the blue palette of the title and abstract.
\newtcolorbox{highlightbox}[1][Key idea]{
  colback=jlublue!4,colframe=jlublue,boxrule=0.5pt,
  arc=1.5mm,title=#1,fonttitle=\bfseries,
  left=2mm,right=2mm,top=2mm,bottom=2mm
}

\usetikzlibrary{positioning}


% Jilin University research paper: replace the title and author information.
\title{Your Paper Title:\\A Clear and Concise Subtitle}
\author[1]{First Author}
\author[1]{Second Author}
\author[1\corrauthor]{Third Author}
\affiliation[1]{School of Your Discipline, Jilin University, Changchun, China}
% For a collaborating institution, add \affiliation[2]{Institution, City, Country}
% and use \author[1,2]{Author Name}.
\contribution[\mathsection]{Corresponding author}
\correspondence{\email{author@example.com}}
\checkdata[Code]{\url{https://example.com/project}}
% The Jilin University emblem is included; use \paperlogo{} to hide it.
\paperlogo{figures/jlu.pdf}
% \date{October 2026}

\abstract{%
State the problem and explain why it matters. Briefly describe the limitation
of existing approaches, introduce your method, and summarize the main results.
Conclude with the practical significance of your findings. Replace this text
with a self-contained abstract of approximately 150--250 words.
}

\begin{document}
\maketitle
\section{Introduction}
Introduce the research problem, its motivation, and the setting in which it
arises. Explain the gap between current methods and the desired outcome.
Use citations to support specific claims, for example the standard reference
on typesetting by \citet{lamport1994latex}.

\paragraph{Contributions.}
Summarize the contributions in concrete terms:
\begin{itemize}
  \item A formulation of the problem and its main challenges.
  \item A method that addresses the identified limitations.
  \item An evaluation with appropriate baselines and ablations.
\end{itemize}

\section{Related Work}
Organize prior work by theme and explain how your approach relates to it.
Multiple references can be cited together~\citep{lamport1994latex,knuth1984texbook}.

\section{Method}
\subsection{Problem Formulation}
Let $\vect{x}_i\in\R^d$ denote an input and $y_i$ its target. Given a dataset
$\mathcal{D}=\{(\vect{x}_i,y_i)\}_{i=1}^N$, we optimize
\begin{equation}
  \mathcal{L}(\theta)=\frac{1}{N}\sum_{i=1}^N
  \ell\bigl(f_\theta(\vect{x}_i),y_i\bigr).
  \label{eq:objective}
\end{equation}
Define all symbols before introducing the method.

\subsection{Proposed Approach}
Describe the steps needed to reproduce the approach. Refer to
\cref{eq:objective} for the objective and \cref{fig:pipeline} for an overview.
\begin{highlightbox}
Replace this box with a central insight, research question, or concise
summary of the proposed method.
\end{highlightbox}

\begin{figure}[t]
  \centering
  \begin{tikzpicture}[node distance=4mm,
    stage/.style={draw=jlublue,fill=jlublue!5,rounded corners,
                  minimum width=20mm,minimum height=10mm,font=\small}]
    \node[stage] (input) {Input};
    \node[stage,right=of input] (method) {Method};
    \node[stage,right=of method] (output) {Output};
    \draw[->,thick,color=jlublue] (input) -- (method);
    \draw[->,thick,color=jlublue] (method) -- (output);
  \end{tikzpicture}
  % To use an external figure instead:
  % \includegraphics[width=\linewidth]{your-figure.pdf}
  \caption{An example pipeline. Replace it with an overview of your method.}
  \label{fig:pipeline}
\end{figure}

\begin{theorem}[Example statement]
If $\ell$ is nonnegative, then $\mathcal{L}(\theta)\geq 0$.
\end{theorem}
\begin{proof}
Every summand in \cref{eq:objective} is nonnegative.
\end{proof}

\section{Experiments}
\subsection{Experimental Setup}
Describe datasets, baselines, metrics, implementation details, and resources.
Report enough information for another researcher to reproduce the evaluation.

\subsection{Main Results}
\Cref{tab:results} illustrates a compact three-line table. All numbers below
are placeholders, rather than measured experimental results.
\begin{table}[t]
  \centering
  \caption{Example results (illustrative values only).}
  \label{tab:results}
  \begin{tabular}{lrr}
    \toprule
    Method & Score $\uparrow$ & Time $\downarrow$ \\
    \midrule
    Baseline A & 72.1 & 1.20 \\
    Baseline B & 75.3 & 1.05 \\
    Ours & \textbf{79.4} & \textbf{0.92} \\
    \bottomrule
  \end{tabular}
\end{table}
Discuss effect sizes and uncertainty, along with relevant failure cases.
Use \texttt{table*} or \texttt{figure*} for content spanning both columns.

\subsection{Ablation Study}
Isolate the effect of each component and explain what the results reveal.

\section{Conclusion}
Summarize the findings and their implications. State the limitations and
identify concrete directions for future work.

\section*{Acknowledgments}
Add acknowledgments and funding information here when appropriate.

\bibliographystyle{plainnat}
\bibliography{references}

\clearpage
\onecolumn
\appendix
\section{Additional Details}
\label{app:details}
The appendix uses a single-column layout for longer derivations, additional
experiments, and implementation details. Remove \verb|\onecolumn| in
\texttt{main.tex} if you prefer a two-column appendix.

\subsection{Reproducibility Checklist}
Record dataset versions, preprocessing, hyperparameters, random seeds,
software versions, and the hardware used for the experiments.

\end{document}
